\subsection{Linear Representations of Algebras}

DEFINITION 1. --- Let M be a K-module. A linear representation of the algebra A in M is a K-homomorphism \(\pi\) from A to the algebra \(\operatorname{End}_{\mathrm{K}}(\mathrm{M})\) .

We also say that the pair \((\mathrm{M},\pi)\) is a linear representation of the algebra A. When M is a free K-module, the dimension of M is called the degree (or dimension) of \(\pi\) .

Let \(\pi\) be a linear representation of A in M. The additive law on M and the external law \((a,x)\mapsto\pi(a)(x)\) define a left A-module structure on M, which we denote by \(M_{\pi}\) . We say that \(M_{\pi}\) is the module of the representation \(\pi\) . The K-module structure on M is obtained from the A-module structure on \(M_{\pi}\) by restriction of scalars.

Conversely, let E be a left A-module. Let M be the K-module deduced from E by restriction of scalars from A to K, and let \(\pi\) be the homomorphism \(a \mapsto a_{E}\) from A to \(\operatorname{End}_{\mathrm{K}}(\mathrm{M})\) . Then \(\pi\) is a representation of A in M, and we have \(E = M_{\pi}\) . We say that \(\pi\) is the linear representation associated with the A-module E.

Studying A-modules or linear representations of A amounts to the same. We will translate certain definitions concerning modules to the language of linear representations.

Let \(\pi\) be a linear representation of A in M. The kernel of the homomorphism \(\pi\) is a two-sided ideal of A, which is simply the annihilator of the A-module \(\mathrm{M}_{\pi}\) . The homomorphism \(\pi\) is injective if and only if the A-module \(\mathrm{M}_{\pi}\) is faithful; we then say that \(\pi\) is a faithful representation of A.

Let \((\mathrm{M},\pi)\) and \((\mathrm{M}^{\prime},\pi^{\prime})\) be linear representations of A. An A-linear mapping from \(\mathrm{M}_{\pi}\) to \(\mathrm{M}_{\pi^{\prime}}^{\prime}\) is a K-linear mapping \(u:\mathrm{M}\to\mathrm{M}^{\prime}\) that satisfies the relation

\[
\pi^ {\prime} (a) \circ u = u \circ \pi (a)\tag{1}
\]

for \(a \in \mathbf{A}\) . An isomorphism from \(\mathrm{M}_{\pi}\) to \(\mathrm{M}_{\pi'}'\) is an isomorphism of K-modules \(u: \mathrm{M} \to \mathrm{M}'\) satisfying the condition

\[
\pi^ {\prime} (a) = u \circ \pi (a) \circ u ^ {- 1}\tag{2}
\]

for \(a \in A\) . We say that \(\pi\) and \(\pi'\) are isomorphic if the A-modules \(M_{\pi}\) and \(M'_{\pi'}\) are isomorphic. We say that \(\pi\) is a subrepresentation (resp. a quotient representation) of \(\pi'\) if \(M_{\pi}\) is a submodule (resp. a quotient module) of \(M'_{\pi'}\) .

Suppose given a family \((\mathrm{M}_{i},\pi_{i})\) of linear representations of A. Let M be the K-module that is the direct sum of the \(M_{i}\) ; for any \(a\in A\) , let \(\pi(a)\) be the endomorphism \((x_{i})_{i\in\mathrm{I}}\mapsto(\pi_{i}(a)(x_{i}))_{i\in\mathrm{I}}\) of M. Then \(\pi\) is a linear representation of A in M. The A-module \(M_{\pi}\) is the direct sum of the A-modules \((\mathrm{M}_{i})_{\pi_{i}}\) ( \(i\in I\) ); we say that \(\pi\) is the direct sum of the \(\pi_{i}\) , and we write \(\pi=\bigoplus\pi_{i}\) .

We say that the representation \(\pi\) of A in M is simple or irreducible if the A-module \(M_{\pi}\) is simple; we say that the representation \(\pi\) of A on M is semisimple or completely reducible if the A-module \(M_{\pi}\) is semisimple.

Let \(\pi\) be a linear representation of A in M. Let \(\mathrm{M}^*\) be the K-module dual to M. The mapping \(a\mapsto{}^{t}(\pi (a))\) from \(\mathrm{A}^{\circ}\) to \(\mathrm{End}_{\mathrm{K}}(\mathrm{M}^{*})\) is the transposed representation of \(\pi\) .

Let L be a commutative K-algebra, and let \((M,\pi)\) be a linear representation of the K-algebra A. The homomorphism \(\pi_{(L)}:A_{(L)}\to\operatorname{End}_{L}(M_{(L)})\) corresponding to the \(A_{(L)}\) -module structure on \(M_{(L)}\) is a linear representation of the L-algebra \(A_{(L)}\) . We say that \(\pi_{(L)}\) is the linear representation of the algebra \(A_{(L)}\) deduced from the representation \(\pi\) by extending the ring of scalars K to L.

Suppose that K is a field and that L is a nonzero commutative K-algebra.

Let \(\pi\) and \(\pi'\) be linear representations of the algebra A. It follows from VIII, p.~37, Theorem 3 that the representations \(\pi\) and \(\pi'\) are isomorphic if and only if \(\pi_{(L)}\) and \(\pi_{(L)}'\) are.

Suppose that K is a field. Let L be an extension of K. Consider the Grothendieck group \(R_{\mathrm{K}}(\mathrm{A})\) (resp. \(R_{\mathrm{L}}(\mathrm{A}_{(\mathrm{L})})\) ) of the A-modules that are finite-dimensional over K (resp. of the \(A_{(L)}\) -modules that are finite-dimensional over L). We have seen that the group homomorphism

\[
u: \mathrm{R} _ {\mathrm{K}} (\mathrm{A}) \longrightarrow \mathrm{R} _ {\mathrm{L}} (\mathrm{A} _ {(\mathrm{L})})
\]

defined by extension of scalars is injective; moreover, an element \(\xi \in \mathrm{R}_{\mathrm{K}}(\mathrm{A})\) is effective if and only if \(u(\xi)\) is (VIII, p.~195, Theorem 1).

